\section{Conclusion and future work}
\label{sec:conclusion}

This project built and compared three complete approaches to predicting loan default on the Credit Risk Dataset.

The \textbf{traditional pipeline} is fully configurable, from imputation to the classifier, and every option was evaluated with controlled cross-validated ablations.
Scaling decides the fate of distance- and margin-based models, engineered features, IQR clipping and a one-hot grade help the linear models, and PCA and univariate feature selection never help reliably on this low-dimensional table.
After tuning, histogram-based gradient boosting is the best model, with a test PR-AUC of \BestPrAuc{}.
At its F2-tuned threshold it detects \BestRecall{} of the defaulters with a precision of \BestPrecision{}.
A wider hyper-parameter search and a stacked ensemble did not improve on it, which indicates that the result is close to what these features allow.

The \textbf{deep learning pipeline} with entity embeddings outperforms every linear model and produces reusable learned features, but it stays below the tree ensembles, in line with the literature on tabular data.

The \textbf{Bayesian network}, implemented from scratch for the talented-engineering extension, is the least accurate classifier but the most transparent model.
TAN is the best structure for prediction, whereas a better model of the joint distribution, as learned by hill climbing on BIC, does not translate into better classification; comparing exact query answers with empirical frequencies exposes the independence assumptions of each structure.
The inference experiments reproduce the textbook behaviour: variable elimination is exact and fast on these small networks, likelihood weighting and Gibbs sampling converge, and rejection sampling collapses once many variables are observed.

\paragraph{Future work.}
\begin{itemize}[nosep]
  \item Calibrate the probabilities (isotonic or Platt scaling) and choose the threshold by minimizing an explicit expected loss.
  \item Explain individual decisions with SHAP values, as lenders must justify rejections.
  \item Audit the model for disparate impact across age groups.
  \item Learn Bayesian-network parameters with expectation maximization over the incomplete rows, and try continuous (conditional Gaussian) nodes to avoid discretization.
  \item Validate on a time-ordered split if dated loan data become available.
\end{itemize}
